\documentclass[11pt,a4paper]{article}
\usepackage[margin=0.85in,top=1in,bottom=1in]{geometry}
\usepackage{amsmath,amssymb,amsfonts}
\usepackage{graphicx}
\usepackage{float}
\usepackage{enumitem}
\usepackage{microtype}
\usepackage{caption}
\usepackage[hidelinks]{hyperref}
\usepackage[most,breakable]{tcolorbox}
\tcbuselibrary{breakable,skins}
\usepackage[utf8]{inputenc}
\usepackage[T1]{fontenc}
\usepackage{lmodern}
\captionsetup{font=small,labelfont=bf,skip=4pt}
\setlist{leftmargin=1.5em,topsep=2pt}

%% Breakable card — prevents content cutoff across pages
\newtcolorbox{solcard}[3]{
  breakable, enhanced,
  colback=white, colframe=gray!50, arc=2pt,
  title={\textbf{#1}\hfill\textsf{\small #3\,pts}},
  fonttitle=\normalsize\sffamily\bfseries,
  before upper={\small\textit{#2}\par\smallskip\hrule height 0.4pt\smallskip},
  top=4pt, bottom=6pt, left=7pt, right=7pt,
  parbox=false,
}

\title{\textbf{Conductors in a magnetic field --- Solution}}
\author{\normalsize X24 (2024) — English translation}
\date{}
\begin{document}\maketitle\vspace{-0.5em}
\begin{solcard}{A1}{Let the moment of time $t_0=0$ the load be at the origin of coordinates, and the projection of its velocity onto the axis $x$ is equal to $v_0$. Determine the dependences of the coordinate $x(t)$ and speed $v_x(t)$ of the load on time $t$. Express your answer using $v_0$, $\gamma$, $\omega_0$ and $t$.}{0.60}

Let's write down the equation of motion of the load: $$m\ddot{x}=-\beta\dot{x}-kx\Rightarrow \ddot{x}+\cfrac{\beta}{m}\dot{x}+\cfrac{k}{m}x=0{.} $$Taking into account introduce notation: $$\ddot{x}+2\gamma\dot{x}+\omega^2_0x=0{.} $$We seek solution in complex form: $$x=Re\left\{Ae^{\lambda t}\right\}{,} $$where $A\neq 0$ and $\lambda$ – some complex number. Then: $$A\left(\lambda^2+2\gamma\lambda+\omega^2_0\right)=0\Rightarrow{\lambda=-\gamma\pm{i\sqrt{\omega^2_0-\gamma^2}}}{.} $$Solution represent itself sum solution, corresponding different value $\lambda$: $$x(t)=e^{-\gamma t}Re\left\{A_1e^{i\sqrt{\omega^2_0-\gamma^2}t}+A_2e^{-i\sqrt{\omega^2_0-\gamma^2}t}\right\}=Ce^{-\gamma t}\sin\left(\sqrt{\omega^2_0-\gamma^2}t+\varphi_0\right){,} $$where $C$ and $\varphi_0$ are real numbers. Differentiate: $$v_x(t)=Ce^{-\gamma t}\left(\sqrt{\omega^2_0-\gamma^2}\cos\left(\sqrt{\omega^2_0-\gamma^2}t+\varphi_0\right)-\gamma\sin\left(\sqrt{\omega^2_0-\gamma^2}t+\varphi_0\right)\right){.} $$For moment time $t=0$ have: $$\begin{cases} x(0)=C\sin\varphi_0\\ v_x(0)=C\left(\sqrt{\omega^2_0-\gamma^2}\cos\varphi_0-\gamma\sin\varphi_0\right) \end{cases} $$Taking into account initial condition: $$\begin{cases} x(0)=0\\ v_x(0)=v_0 \end{cases} \Rightarrow \begin{cases} \varphi_0=0\\ C=\cfrac{v_0}{\sqrt{\omega^2_0-\gamma^2}} \end{cases} $$Final we get:

\par\smallskip\noindent\textbf{Answer:} $$x(t)=\cfrac{v_0}{\sqrt{\omega^2_0-\gamma^2}}e^{-\gamma t}\sin\left(\sqrt{\omega^2_0-\gamma^2}t\right){.} $$
\end{solcard}

\begin{solcard}{A2}{Obtain an exact expression for $Q$. Express your answer in terms of $\omega_0$ and $\gamma$.}{0.40}

Let's rewrite the expression for the quality factor in the following form: $$Q=\cfrac{2\pi}{1-\left(\cfrac{v_1}{v_0}\right)^2}{,} $$where $v_1$ – the projection of the weight velocity at the second passage of the origin with the more same direction velocity. Quantity velocity $v_1$ be achieved through period $T$, equal: $$T=\cfrac{2\pi}{\sqrt{\omega^2_0-\gamma^2}}{.} $$Then: $$\cfrac{v_1}{v_0}=e^{-\gamma T}=e^{-2\pi\gamma/\sqrt{\omega^2_0-\gamma^2}}{,} $$from where we get:

\par\smallskip\noindent\textbf{Answer:} $$Q=\cfrac{2\pi}{1-e^{-4\pi\gamma/\sqrt{\omega^2_0-\gamma^2}}}{.} $$
\end{solcard}

\begin{solcard}{A3}{Derive an approximate expression for the quality factor $Q$ at weak attenuation ($\gamma\ll\omega_0$). Express your answer using $m$, $k$ and $\beta$.}{0.20}

The expansion of the denominator at $\gamma\ll{\omega_0}$ is as follows: $$1-e^{-4\pi\gamma/\sqrt{\omega^2_0-\gamma^2}}\approx 1-e^{-4\pi\gamma/\omega_0}\approx 1-\left(1-\cfrac{4\pi\gamma}{\omega_0}\right)=\cfrac{4\pi\gamma}{\omega_0}{,} $$from which obtain approximation for the quality factor for weak damping $Q$: $$Q\approx\cfrac{\omega_0}{2\gamma}{.} $$Substitute $\omega_0$ and $\gamma$, we obtain:

\par\smallskip\noindent\textbf{Answer:} $$Q\approx\cfrac{\sqrt{mk}}{\beta(H)}{.} $$
\end{solcard}

\begin{solcard}{B1}{The deviation $x$ of the load from the position depends on time $t$ as follows: $$x(t)=A\sin\left(\Omega t+\varphi_0\right) $$Find $A$ and $\varphi_0$. Answer express through $A_0$, $\Omega$, $\omega_0$ and $\gamma$.}{0.60}

The equation of motion is as follows: $$m\ddot{x}=k(A_0\sin\Omega t-x)-\beta\dot{x}{,} $$from which: $$\ddot{x}+2\gamma\dot{x}+\omega^2_0\Delta{x}=\omega^2_0A_0\sin\Omega t{.} $$We seek solution in form: $$x=Re\left\{\hat{A}e^{i\Omega t}\right\}{,} $$where $\hat{A}\neq 0$ – some complex number. Then: $$\hat{A}\left((\omega^2_0-\Omega^2)+2i\Omega\gamma\right)=\omega^2_0A_0e^{-i\pi/2}{.} $$Express $A$: $$\hat{A}=\cfrac{\omega^2_0A_0\left((\omega^2_0-\Omega^2)-2i\Omega\gamma\right)e^{-i\pi/2}}{\left((\omega^2_0-\Omega^2)^2+4\gamma^2\Omega^2\right)}=\cfrac{A_0\omega^2_0e^{-i(\pi/2-\varphi)}}{\sqrt{(\omega^2_0-\Omega^2)^2+4\gamma^2\omega^2_0}}{,} $$where $\varphi$ – argument complex number $z=(\omega^2_0-\Omega^2) -2i\Omega\gamma$. Extracting the real part, we find: $$\Delta{x}(t)=\cfrac{A_0\omega^2_0}{\sqrt{(\omega^2_0-\Omega^2)^2+4\gamma^2\Omega^2}}\sin\left(\Omega t+\varphi\right){.} $$Thus: $$A=\cfrac{A_0\omega^2_0}{\sqrt{(\omega^2_0-\Omega^2)^2+4\gamma^2\Omega^2}}\qquad \varphi_0=\varphi{.} $$For determining $\varphi$ there are three cases: $$\varphi=\begin{cases} -\arctan\cfrac{2\gamma\Omega}{\omega^2_0-\Omega^2}\quad\text{for}\quad \Omega{<}\omega_0\\ -\cfrac{\pi}{2}\quad\text{for}\quad \Omega=\omega_0\\ -\pi-\arctan\cfrac{2\gamma\Omega}{\omega^2_0-\Omega^2}\quad\text{for}\quad \Omega{>}\omega_0 \end{cases} $$So way:

\par\smallskip\noindent\textbf{Answer:} $$A=\cfrac{A_0\omega^2_0}{\sqrt{(\omega^2_0-\Omega^2)^2+4\gamma^2\Omega^2}}{.} $$$$\varphi_0=\begin{cases} -\arctan\cfrac{2\gamma\Omega}{\omega^2_0-\Omega^2}\quad\text{for}\quad \Omega{<}\omega_0\\ -\cfrac{\pi}{2}\quad\text{for}\quad \Omega=\omega_0\\ -\pi-\arctan\cfrac{2\gamma\Omega}{\omega^2_0-\Omega^2}\quad\text{for}\quad \Omega{>}\omega_0 \end{cases} $$
\end{solcard}

\begin{solcard}{B2}{Obtain exact expressions for the resonant cyclic frequency $\Omega_{res}$ and the corresponding oscillation amplitude $A_{res}$. Express your answers using $\omega_0$, $\gamma$ and $A_0$. Consider that $\gamma\sqrt{2}<\omega_0$.}{0.30}

Differentiating the denominator with respect to $\Omega^2$, we obtain: $$-2(\omega^2_0-\Omega^2)+4\gamma^2=0{,} $$from where:


Substituting $\Omega_\text{res}$ into the expression for $A$, we find:

\par\smallskip\noindent\textbf{Answer:} $$\Omega_\text{res}=\sqrt{\omega^2_0-2\gamma^2}{.} $$
\end{solcard}

\begin{solcard}{B3}{Derive approximate expressions for $\Omega_{res}$, $A_{res}$ and $\Delta{\omega}$ under weak damping ($\gamma\ll{\omega_0}$). Express your answers using $A_0$, $\omega_0$ and $\gamma$.}{0.30}

Simplified expressions for $\Omega_\text{res}$ and $A_\text{res}$ are obtained trivially and take the following form:


Consider the cyclic frequency $\Omega=\Omega_\text{res}+\Delta\Omega$, where $\Delta\Omega\ll\Omega_\text{res}$. The value of $\Omega_\text{res}$ can be considered equal to $\omega_0$, since: $$\Omega_\text{res}=\sqrt{\omega^2_0-2\gamma^2}\approx \omega_0-\cfrac{\gamma^2}{\omega_0}{.} $$Deviation $\Omega_\text{res}$ from $\omega_0$ represent itself quantity second order smallness. Then for under the radical expression obtain: $$\left(\omega^2_0-\Omega^2\right)^2+4\gamma^2\Omega^2\approx \left(\omega^2_0-\left(\omega_0+\Delta\Omega\right)\right)^2+4\gamma^2\left(\omega_0+\Delta\Omega\right)^2\approx 4\omega^2_0\Delta\Omega^2+4\gamma^2\omega^2_0{.} $$The quantity $\Delta\Omega$ is such that under the radical the expression is twice larger corresponding resonance. Hence: $$4\Delta{\Omega}^2\omega^2_0+4\gamma^2\omega^2_0=8\gamma^2\omega^2_0\Rightarrow \Delta{\Omega}_{1{,}2}=\pm\gamma{.} $$Since $\Delta\omega=\Delta\Omega_1-\Delta\Omega_2$, we get:

\par\smallskip\noindent\textbf{Answer:} $$\Omega_\text{res}\approx{\omega_0}\qquad A_\text{res}\approx{\cfrac{A_0\omega_0}{2\gamma}}{.} $$
\end{solcard}

\begin{solcard}{C1}{Find the induction $B_x$ of the magnetic field of the ring on its axis at the point with coordinate $x$. Express your answer in terms of $x$, $R$, $I$ and the magnetic constant $\mu_0$.}{0.30}

From the Biot-Savart-Laplace law: $$d\vec{B}=\cfrac{\mu_0}{4\pi}\cfrac{\left[\vec{r}\times d\vec{r}\right]}{r^3}{,} $$where $\vec{r}$ – radius–vector element ring relatively point with coordinate $x$. For each point ring: $$r=\sqrt{R^2+x^2}{.} $$Integral by circuit, from the vector product we obtain from its geometric meaning: $$\oint_L\left[\vec{r}\times d\vec{r}\right]=2\vec{S}=2\pi{R}^2\vec{e}_x{,} $$from:

\par\smallskip\noindent\textbf{Answer:} $$B_x(x)=\cfrac{\mu_0IR^2}{2\left(R^2+x^2\right)^{3/2}}{.} $$
\end{solcard}

\begin{solcard}{C2}{Determine the magnetic moment $\vec{m}$ of the disk. Express your answer in terms of $\vec{e}_x$, $r_0$, $h$, $\rho$ and $\dot{B}$.}{1.00}

From Faraday's law of electromagnetic induction we obtain: $$\mathcal{E}_\text{ind}=-\cfrac{d\Phi}{dt}=-\pi{r}^2\dot{B}=2\pi rE_\text{vort}{,} $$from which: $$E_\text{vort}=-\cfrac{\dot{B}r}{2}{.} $$From Ohm's law in differential form: $$\vec{j}=\cfrac{\vec{E}}{\rho}{,} $$from which: $$j(r)=-\cfrac{\dot{B}r}{2\rho}{.} $$The moment of the ring current of height $h$ and thickness $dr$ equal: $$dm_x=\pi{r}^2dI=\pi{r}^2jhdr=-\cfrac{\pi\dot{B}h}{2\rho}r^3dr{,} $$from which: $$m_x=-\cfrac{\pi\dot{B}h}{2\rho}\int\limits_0^{r_0}r^3dr{.} $$Integrating, we find:

\par\smallskip\noindent\textbf{Answer:} $$\vec{m}=-\vec{e}_x\cdot{\cfrac{\pi{r}^4_0h\dot{B}}{8\rho}}{.} $$
\end{solcard}

\begin{solcard}{C3}{Determine the magnetic moment $\vec{m}$ of the ball. Express your answer in terms of $\vec{e}_x$, $R_0$, $\rho$ and $\dot{B}$.}{0.50}

Let's use spherical coordinates (that is, we will count the angle $\theta$ from the positive direction of the axis $x$). Then for $r_0$ and $h$ we have: $$r_0=R_0\sin\theta\qquad h=d(R_0(1-\cos\theta))=R_0\sin\theta d\theta{.} $$Then for element magnetic moment sphere have: $$dm_x=-\cfrac{\pi{R}^5_0\dot{B}}{8\rho}\sin^5\theta d\theta{,} $$from which: $$m_x=-\cfrac{\pi{R}^5_0\dot{B}}{8\rho}\int\limits_{0}^{\pi}\sin^5\theta d\theta{.} $$Integrate obtain expression with means substitution $x=\cos\theta$: $$\int\limits_{0}^{\pi}\sin^5\theta d\theta=\int\limits_{-1}^{1}(1-x^2)^2dx=\int\limits_{-1}^{1}(1-2x^2+x^4)dx=\left(x-\cfrac{2x^3}{3}+\cfrac{x^5}{5}\right)\Biggl|_{-1}^{1}=\cfrac{16}{15}{,} $$from:

\par\smallskip\noindent\textbf{Answer:} $$\vec{m}=-\vec{e}_x\cdot\cfrac{2\pi{R}^5_0\dot{B}}{15\rho}{.} $$
\end{solcard}

\begin{solcard}{C4}{Obtain the time derivative of the magnetic field induction of the ring at the center of the ball $dB_x/dt$, which is equivalent to the value $\dot{B}$. Express your answer in terms of $v$, $I$, $R$, $x$ and the magnetic constant $\mu_0$.}{0.40}

When the ball moves, the magnetic field induction $B_x(x)=B_x(x(t))$, i.e., is a complex function of time, so we have: $$\dot{B_x}=\cfrac{dB_x}{dt}=\cfrac{dB_x}{dx}\cfrac{dx}{dt}=v\cfrac{dB_x}{dx} $$Find derivative $dB_x/dx$: $$\cfrac{dB_x}{dx}=-\cfrac{3\mu_0IR^2x}{2(R^2+x^2)^{\frac{5}{2}}}{,} $$from where:

\par\smallskip\noindent\textbf{Answer:} $$\dot{B}=-\cfrac{3\mu_0IR^2xv}{2(R^2+x^2)^{\frac{5}{2}}}{.} $$
\end{solcard}

\begin{solcard}{C5}{Find the proportionality factor $\beta(x)$. Express your answer in terms of $I$, $R$, $x$, $R_0$, $\rho$ and the magnetic constant $\mu_0$.}{0.50}

Since the ball moves along the $x$ axis: $$\vec{F}=m_x\cdot\cfrac{\partial\vec{B}}{\partial x}=\vec{e}_x\cdot m_x\cfrac{dB_x}{dx}{.} $$For magnetic moment $m_x$ have: $$m_x=-\cfrac{2\pi{R}^5_0v}{15\rho}\cdot\cfrac{dB_x}{dx}{,} $$from which: $$F_x=-\cfrac{2\pi{R}^5_0}{15\rho}\left(\cfrac{dB_x}{dx}\right)^2v{.} $$Substitute $dB_x/dx$, we find:

\par\smallskip\noindent\textbf{Answer:} $$\beta(x)=\cfrac{3\pi\mu^2_0I^2R^4R^5_0x^2}{10\rho(R^2+x^2)^5}{.} $$
\end{solcard}

\begin{solcard}{C6}{Determine the resistivity $\rho$ of the ball used in the first experiment. Express your answer in terms of $m$, $k$, $R_0$, $R$, $H$, $I$ and the magnetic constant $\mu_0$.}{0.80}

The first peak is reached at $\Delta{x}\approx{47~\text{a.u.}}$, and the 12th at $\Delta{x}\approx{6~\text{a.u.}}$. Hence: $$\cfrac{\gamma}{\sqrt{\omega^2_0-\gamma^2}}\approx{\cfrac{\ln\left(\cfrac{47}{6}\right)}{11\cdot{2\pi}}}\approx{0{,}0297}\Rightarrow \cfrac{\gamma}{\omega_0}\approx 0{.}03\ll{1}{.} $$At the same time have: $$\cfrac{\omega_0}{2\gamma}=\cfrac{\sqrt{mk}}{\beta(H)}{,} $$from which: $$\beta(H)\approx{\cfrac{2\gamma\sqrt{mk}}{\omega_0}}{,} $$and finally

\par\smallskip\noindent\textbf{Answer:} $$\rho=15{.}7\cdot\cfrac{\mu^2_0I^2R^5_0R^4H^2}{\sqrt{mk}(R^2+H^2)^5} $$
\end{solcard}

\begin{solcard}{C7}{Determine the resistivity $\rho$ of the ball used in the second experiment. Express your answer in terms of $m$, $k$, $R_0$, $R$, $H$, $I$ and the magnetic constant $\mu_0$.}{0.70}

From the expression for the resonant amplitude we find: $$Q=\cfrac{A_\text{res}}{A_0}=\cfrac{\sqrt{mk}}{\beta(H)}\approx 25 $$from where finally:

\par\smallskip\noindent\textbf{Answer:} $$\rho=23{.}6\cfrac{\mu^2_0I^2R^5_0R^4H^2}{\sqrt{mk}(R^2+H^2)^5} $$
\end{solcard}

\begin{solcard}{D1}{Determine the induction $B_z$ of the magnetic field of the solenoid, as well as its derivative $dB_z/dz$ at the point with coordinate $z$. Express your answer using $\mu_0$, $n$, $I$, $R$ and $z$.}{0.60}

From the solid angle theorem for the magnetic field we have: $$B_z=\cfrac{\mu_0i\Omega_\text{side}}{4\pi}{,} $$where $i=In$ – linear density current solenoid, and $\Omega_\text{side}$ – the solid angle subtended by its lateral surface. The solid angle subtended by the lateral surface of the solenoid equals the solid angle that can be seen it base from point with coordinate $z$, therefore have: $$\Omega_\text{side}=2\pi(1-\cos\alpha)=2\pi\left(1-\cfrac{z}{\sqrt{z^2+R^2}}\right){.} $$Thus:


Differentiating, we find:

\par\smallskip\noindent\textbf{Answer:} $$B_z=\cfrac{\mu_0nI}{2}\left(1-\cfrac{z}{\sqrt{z^2+R^2}}\right){.} $$
\end{solcard}

\begin{solcard}{D2}{Determine the linear current density $i$ on the surface of the cylinder at the point with coordinate $z$. Express your answer using $\mu_0$, $x$ and $dB_z(z)/dz$.}{1.00}

Since the magnetic field induction outside the cylinder can be considered equal to the magnetic field induction of the solenoid, we have: $$B_{z(in)}=B_z(z-x)\qquad B_{z(out)}=B(z){.} $$From the circulation theorem for the magnetic-field induction we obtain: $$(B_{z(in)}-B_{z(out)})\approx -x\cfrac{dB_z}{dz}=\mu_0ix{,} $$from where:

\par\smallskip\noindent\textbf{Answer:} $$i(z)=-\cfrac{x}{\mu_0}\cfrac{dB_z}{dz}{.} $$
\end{solcard}

\begin{solcard}{D3}{Determine the force $F_x$ acting on the cylinder from the magnetic field of the solenoid. Express your answer in terms of $\mu_0$, $r$, $R$, $n$, $I$ and $x$.}{1.50}

The magnetic moment of an infinitesimal section of a cylinder is: $$dm_z=i(z)\pi r^2dz{.} $$For act on it force $dF_x$ have: $$dF_x=dm_z\cfrac{dB_z}{dz}{.} $$Substitute expression for $i$, obtain: $$dF_x=-\cfrac{\pi{r}^2x}{\mu_0}\left(\cfrac{dB_z}{dz}\right)^2dz{.} $$Use expression for $dB_z/dz$: $$dF_x=-\cfrac{\mu_0\pi{r}^2n^2I^2xR^4dz}{4(R^2+z^2)^3}{,} $$from which taking into account the larger distance of the end of the cylinder from the base of the solenoid: $$F_x\approx -\cfrac{\mu_0\pi{r}^2n^2I^2R^4x}{4}\int\limits_{-\infty}^{\infty}\cfrac{dz}{(R^2+z^2)^3}{.} $$Using the substitution $z = R\tan\varphi$ we obtain: $$F_x=\cfrac{\mu_0\pi{r}^2n^2I^2R^4x}{4}\int\limits_{-\infty}^{\infty}\cfrac{dz}{(R^2+z^2)^3}=\cfrac{\mu_0\pi{r}^2n^2I^2x}{4R}\int\limits_{-\pi/2}^{\pi/2}\cos^4\varphi d\varphi{.} $$Let us compute last integral: $$\int\limits_{-\pi/2}^{\pi/2}\cos^4\varphi d\varphi=\int\limits_{-\pi/2}^{\pi/2}\left(\cfrac{1+\cos 2\varphi}{2}\right)^2d\varphi=\int\limits_{-\pi/2}^{\pi/2}\left(\cfrac{1}{4}+\cfrac{\cos 2\varphi}{2}+\cfrac{\cos^22\varphi}{4}\right)d\varphi= $$ $$=\int\limits_{-\pi/2}^{\pi/2}\left(\cfrac{1}{4}+\cfrac{\cos 2\varphi}{2}+\cfrac{1}{4}\left(\cfrac{1+\cos 4\varphi}{2}\right)\right)d\varphi=\int\limits_{-\pi/2}^{\pi/2}\left(\cfrac{3}{8}+\cfrac{\cos 2\varphi}{2}+\cfrac{\cos 4\varphi}{8}\right)d\varphi= $$ $$=\left(\cfrac{3\varphi}{8}+\cfrac{\sin 2\varphi}{4}+\cfrac{\sin 4\varphi}{32}\right)\Biggl|_{-\pi/2}^{\pi/2}=\cfrac{3\pi}{8}{.} $$So way:

\par\smallskip\noindent\textbf{Answer:} $$F_x=-\cfrac{3\pi^2\mu_0\pi r^2n^2I^2}{32R}x{.} $$
\end{solcard}

\begin{solcard}{D4}{Obtain the dependence of the displacement of the rod $x$ on time $t$. Express your answer in terms of $\mu_0$, $r$, $R$, $n$, $I$ and $m$.}{0.30}

Let us write the equation of motion of the cylinder: $$m\ddot{x}=-\cfrac{3\mu_0\pi^2 r^2n^2I^2}{32R}x\Rightarrow \omega^2_0=\cfrac{3\mu_0\pi^2 r^2n^2I^2}{32mR}\Rightarrow x(t)=A\sin\omega_0t+B\cos\omega_0t{.} $$Determinable constant $A$ and $B$ from initial condition: $$ \begin{cases} x(0)=B=0\\ v_x(0)=\omega_0A=v_0 \end{cases} \Rightarrow A=\cfrac{v_0}{\omega_0} $$Thus:

\par\smallskip\noindent\textbf{Answer:} $$x(t)=v_0\sqrt{\cfrac{32mR}{3\mu_0\pi^2{r}^2n^2I^2}}\sin\sqrt{\cfrac{3\mu_0\pi^2 r^2n^2I^2}{32mR}}t{.} $$
\end{solcard}

\end{document}